\begin{figure}[h]
    \vspace{-3em}
    \centering
    \includegraphics[width=1\linewidth]{images/teaser.pdf}
    \caption{
        10M에서 1B 파라미터까지 확장되는 단일 아키텍처가 과제별 모델을 일관되게 능가한다.
    }
    \label{fig:teaser_results}
\end{figure}


\section{서론}
\label{sec:introduction}
파운데이션 모델은 넓은 데이터 분포 위에서 대규모로 학습된 뒤 다양한 다운스트림 과제로 적응되는 범용 모델이다~\citep{bommasani2021foundation}.
이런 모델은 자연어 처리~\citep{devlin2019bert,brown2020language}와 컴퓨터 비전~\citep{kirillov2023segment,caron2021emerging}을 크게 바꾸어 놓았지만, 멀티 소스 뱅킹 사용자 이력에 대한 적용은 상대적으로 거의 탐구되지 않았다.
현대의 은행과 핀테크는 카드 결제·송금 거래, 상품 사용, 인앱 내비게이션, 고객 커뮤니케이션 같은 이벤트 스트림과 더불어, 계정 보유 기간이나 플랜 같은 정적·일반화된 프로필 상태까지 방대한 양의 데이터를 축적한다.
이런 이벤트 스트림에는 리스크 관리, 상품 분석, 운영에 중요한 신호가 담겨 있지만, 기성 언어모델 토큰화와 아키텍처로는 이를 효율적으로 모델링하기 어렵다.
구조화된 레코드를 텍스트로 직렬화해서 표준 트랜스포머에 넣는 방식은 하나의 베이스라인이 될 수는 있지만, 각 필드 이름과 구분자가 여러 서브워드 토큰으로 쪼개지기 때문에 시퀀스 길이를 크게 부풀린다.
또한 수치값은 숫자 조각으로 분해되면서 크기와 순서 정보가 사라지는데, 이 둘은 금융 추론에서 핵심적이다.
이런 한계 때문에 단순 텍스트 직렬화는 길고 이질적인 뱅킹 사용자 이력에는 실용적이지 않다.

멀티 소스 뱅킹 사용자 이력은 텍스트와 세 가지 점에서 다르다.
첫째, 각 이벤트는 범주형, 수치형, 자유 텍스트 필드가 섞인 가변 길이 레코드다.
둘째, 이력 길이는 롱테일 분포를 보이고, 시간적으로 불규칙하며, 강한 일간 및 주간 주기를 가진다.
셋째, 실제 배포는 엄격한 개인정보 및 규제 제약 아래 이뤄지므로, 무엇을 보고할 수 있는지와 어떤 피처를 특정 결정에 사용할 수 있는지가 제한된다.
이 세 문제를 동시에 처리하는 기성 아키텍처가 없기 때문에, 실무자는 보통 광범위한 피처 엔지니어링이 들어간 과제별 파이프라인을 구축하게 되고, 그 결과 도메인과 상품 사이에서 통계적 강점을 공유하기가 어렵다.

기존 연구는 이 문제의 일부 단면만 다뤘다.
TabTransformer와 FT-Transformer~\citep{huang2020tabtransformer,gorishniy2021revisiting} 같은 표 형식(tabular) 트랜스포머는 고정 스키마 행을 모델링하고, SASRec와 BERT4Rec~\citep{kang2018self,sun2019bert4rec} 같은 순차 추천 모델은 아이템형 상호작용 이력을 다룬다.
금융 파운데이션 모델은 대체로 텍스트나 일반 시계열 토큰화에 집중해 왔으며~\citep{yang2020finbert,wu2023bloomberggpt,yang2023fingpt,jin2024timellm,ansari2024chronos}, nuFormer와 TransactionGPT~\citep{braithwaite2025your,dou2025transactiongpt} 같은 최근 거래 원장(transaction-ledger) 모델은 본 논문의 설정에 더 가깝다.
그러나 이런 모델들은 보통 단일 이벤트 소스만 입력받고 정적 프로필 상태를 생략하며, 좁은 과제 집합에서만 평가된다. 즉, nuFormer는 상품 추천을, TransactionGPT는 이상 탐지와 궤적 생성에 집중한다.
명시적 프로필 상태를 포함하면서 폭넓은 판별적 뱅킹 과제로 전이되는 멀티 소스 인코더 백본은 여전히 문헌에서 비어 있는 자리다.

이 논문에서 저자들은 멀티 소스 뱅킹 사용자 이력을 위한 인코더형 파운데이션 모델 계열인 PRAGMA를 제시한다.
PRAGMA는 멀티 소스 이벤트와 정적 프로필 상태를 함께 담은 대규모 사용자 이력 코퍼스 위에서 마스킹 모델링으로 사전학습된다~(\S\ref{sec:data}).
이질성을 처리하기 위해 수치형·범주형·텍스트형 필드에 맞는 타입별 값 인코딩을 포함한 키--값--시간 토큰화 방식을 적용한다~(\S\ref{sec:tokenisation}).
이렇게 얻어지는 백본은 프로필 상태와 이벤트를 각각 처리하는 두 개의 인코더 가지를 두고, 그 출력을 히스토리 인코더가 융합하는 구조다~(\S\ref{sec:architecture}).

저자들이 인코더 전용·양방향 설계를 택한 이유는, 목표가 개방형 생성이 아니라 판별적 금융 과제를 위한 전이 가능한 표현이기 때문이다.
마스킹 모델링은 각 토큰이 과거와 미래 맥락 모두에 어텐션을 줄 수 있게 해 주며~\citep{devlin2019bert}, 부분적으로만 관측된 이벤트 레코드를 복원하고 전체 이력으로부터 레코드 수준 표현을 학습하는 데 특히 유용하다.
사전학습 이후 PRAGMA는 두 가지 상보적인 방식으로 효율적으로 적응될 수 있다~(\S\ref{sec:protocol}).
\emph{임베딩 프로브(embedding probe)} 설정에서는 백본을 고정한 채, 추출된 임베딩 위에 가벼운 헤드만 학습한다.
\emph{LoRA 파인튜닝} 설정에서는 Low-Rank Adaptation(LoRA)~\citep{hu2022lora}을 적용해 소수의 파라미터만 갱신함으로써, 대부분의 백본을 과제 간 공유하면서도 빠른 특화를 가능하게 한다.

저자들은 신용 평가, 사기 탐지, 커뮤니케이션 인게이지먼트, 정기 거래 탐지, 라이프타임 밸류 예측 등 폭넓은 도메인을 포함하는 내부 다운스트림 벤치마크 위에서 PRAGMA를 평가한다~(\S\ref{sec:downstreams}).
평가된 도메인 전반에서, PRAGMA는 강한 과제별 베이스라인을 일관되게 능가하면서 수작업 피처의 필요성을 줄인다(그림~\ref{fig:teaser_results}).
또한 시퀀스 패킹과 동적 배칭을 포함해, 길고 변동성이 큰 사용자 이력 위에서 PRAGMA를 효율적으로 학습시키기 위한 엔지니어링 선택들도 함께 기술한다~(\S\ref{sec:train_infra}).


저자들의 기여는 다음과 같다.
\begin{itemize}
    \item 멀티 소스 뱅킹 사용자 이력을 위한 인코더형 파운데이션 모델 계열 PRAGMA를 제안한다. 10\,M부터 1\,B 파라미터까지 확장되며, 저자들이 아는 한 소비자 뱅킹 이벤트 시퀀스용으로 공개된 인코더 백본 중 가장 크다. 본 아키텍처는 키--값--시간 토큰화 방식과, 이질적 금융 레코드를 위한 두 가지(branch) 설계를 결합한다. 즉, 프로필 상태 인코더와 이벤트 인코더가 각각 입력을 처리하고 그 결과를 히스토리 인코더가 받는다.
    \item 마스킹 모델링, 시퀀스 패킹, 동적 배칭에 기반해 길고 불규칙한 뱅킹 사용자 이력을 위한 효율적 사전학습 레시피를 기술하고, 사전학습 백본 위에서 LoRA 파인튜닝을 하는 것이 처음부터의 전체 학습과 일관되게 맞먹거나 그 이상임을 보인다.
    \item 단일 사전학습 백본을 여섯 개의 다양한 다운스트림 과제(신용 평가, 사기 탐지, 라이프타임 밸류, 커뮤니케이션 인게이지먼트, 정기 거래 탐지, 상품 추천)에 걸쳐 평가한다. 이는 보통 한두 과제에 한정되었던 기존 거래 원장 모델보다 훨씬 넓은 과제 범위이며, PRAGMA는 강한 과제별 베이스라인을 일관되게 능가하면서 수작업 피처의 필요성을 줄인다.
\end{itemize}
